\subsection{SUBSPACES OF A TOPOLOGICAL SPACE}

Let A be a subset of a topological space X. We have defined the topology induced on A by the topology of X as the inverse image of the latter under the canonical injection \(A \rightarrow X\) (§ 2, no. 3, Example 1). An equivalent definition is as follows:

DEFINITION 1. Let A be a subset of a topological space X. The topology induced on A by the topology of X is that in which the open sets are the intersections with A of open sets of X. The set A with this topology is called a subspace of X.

Example. The topology induced on the set Z of rational integers by the topology of the rational line is the discrete topology, for the intersection of Z and the open interval \(]n - 1/2, n + 1/2[\) is the set \(\{n\}\) .

By Proposition 5 of § 2, no. 3 (or directly from Definition 1) we see that, if B ⊂ A ⊂ X, the subspace B of X is identical with the subspace B of the subspace A of X (transitivity of induced topologies). If S is a subbase (resp. a base) of the topology of X (§ 2, no. 3, Example 3) its trace \(\mathfrak{S}_{\mathrm{A}}\) on A is a subbase (resp. a base) of the topology induced on A.

In all questions which involve the elements or subsets of A, it is essential to distinguish carefully between their properties as points (resp. subsets) of X, and their properties as points (resp. subsets) of the subspace A. We shall make this distinction by using the phrases ``in A'', ``with respect to A'', or ``relative to A'' to refer to properties in the latter category (possibly contrasting them with the phrases ``in X'', ``with respect to X'', ``relative to X'').

An open set of the subspace A need not be open in X; in order that every set which is open in A should be open in X it is necessary and sufficient that A is open in X. The condition is necessary, since A is open in A, and it is sufficient by virtue of \((\mathrm{O}_{\mathrm{II}})\) and Definition 1.

The sets which are closed in A are the intersections with A of the closed sets in X (§ 2, no. 3, Example 1); as above we see that every set which is closed in A is closed in X if and only if A is closed in X.

The neighbourhoods of a point \(x \in A\) relative to A are the intersections with A of neighbourhoods of x relative to X. Every neighbourhood of x relative to A is a neighbourhood of x relative to X if and only if A is a neighbourhood of x in X.

PROPOSITION 1. If A and B are two subsets of a topological space X, and \(B \subset A\) , then the closure of B in the subspace A is the intersection with A of the closure \(\overline{B}\) of B in X.

If \(x \in A\) , every neighbourhood of x in A is of the form \(V \cap A\) , where V is a neighbourhood of x in X. Since \(V \cap B = (V \cap A) \cap B\) , it follows that x lies in the closure of B with respect to A if and only if x lies in the closure of B with respect to X.

COROLLARY. A subset B of A is dense in A if and only if \(\overline{B} = \overline{A}\) in X (i.e.~if and only if \(A \subset \overline{B}\) ).

It follows that if A, B, C are three subsets of X such that \(A \supset B \supset C\) , and if B is dense in A, and C is dense in B, then C is dense in A (transitivity of density), for we have \(\overline{A} = \overline{B} = \overline{C}\) in X.

PROPOSITION 2. Let A be a dense subset of a topological space X; then for each \(x \in \mathbf{A}\) and each neighbourhood V of \(x\) relative to A, the closure \(\overline{\mathbf{V}}\) of V in X is a neighbourhood of \(x\) in X.

For V contains U ∩ A, where U is an open subset of X which contains x, hence \(\overline{V}\) contains U ∩ \(\overline{A} = U\) (§ 1, no. 6, Proposition 5).

PROPOSITION 3. Let \((\mathbf{A}_i)_{i\in I}\) be a family of subsets of a topological space \(\mathbf{X}\) , such that one of the following properties holds:

\begin{enumerate}
\def\labelenumi{\alph{enumi})}
\item
  The interiors of the \(A_{i}\) cover X.
\item
  \((\mathbf{A}_{\iota})_{\iota \in \mathbf{I}}\) is a locally finite closed covering of \(\mathbf{X}\) ( \(\S 1\) , no. 5).
\end{enumerate}

Under these conditions, a subset B of X is open (resp. closed) in X if and only if each of the sets \(B \cap A_{t}\) is open (resp. closed) in \(A_{t}\) .

Clearly if B is open (resp. closed) in X, then \(B \cap A_{t}\) is open (resp. closed) in \(A_{t}\) . Conversely, suppose first that condition a) is satisfied; since \((\mathbb{C}\mathbf{B}) \cap \mathbf{A}_{t} = \mathbf{A}_{t} - (\mathbf{B} \cap \mathbf{A}_{t})\) , it is enough, by duality, to consider the case in which each of the \(B \cap A_{t}\) is open with respect to \(A_{t}\) . In this case \(B \cap \mathring{A}_{t}\) is open in \(\mathring{A}_{t}\) for each \(t \in I\) , and therefore open in X; and since \(B = \bigcup_{i} (B \cap \mathring{A}_{t})\) by hypothesis, it follows that B is open in X.

Now suppose that b) is satisfied; by duality again, we need only consider the case in which each of the \(B \cap A_{t}\) is closed in \(A_{t}\) , and therefore closed in X. Since the family \((B \cap A_{t})\) is locally finite and \(\mathbf{B} = \bigcup_{t} (B \cap A_{t})\) , it follows from §1, no. 5, Proposition 4 that B is closed in X.

Remark. Let \((\mathbf{U}_i)_{i\in I}\) be an open covering of a topological space \(\mathbf{X}\) , and for each \(i\in I\) let \(\mathfrak{B}_i\) be a base of the topology of the subspace \(\mathbf{U}_i\) of \(\mathbf{X}\) ; then it is clear that \(\mathfrak{B} = \bigcup_{i\in I}\mathfrak{B}_i\) is a base of the topology of \(\mathbf{X}\) .

